% =====================================================================
%  Пользовательские команды и сокращения
% =====================================================================

% Краткие ссылки
\newcommand{\figref}[1]{рисунок~\ref{#1}}
\newcommand{\Figref}[1]{Рисунок~\ref{#1}}
\newcommand{\tabref}[1]{таблица~\ref{#1}}
\newcommand{\Tabref}[1]{Таблица~\ref{#1}}
\newcommand{\eqnref}[1]{(\ref{#1})}
\newcommand{\secref}[1]{раздел~\ref{#1}}

% Английские сокращения как термины
\newcommand{\eng}[1]{\textit{#1}}

% Метки общих обозначений
\DeclareMathOperator*{\argmax}{arg\,max}
\DeclareMathOperator*{\argmin}{arg\,min}
\DeclareMathOperator{\softmax}{softmax}
\DeclareMathOperator{\sigm}{\sigma}
\DeclareMathOperator{\KL}{KL}
\DeclareMathOperator{\MMD}{MMD}
\DeclareMathOperator{\PSI}{PSI}

% Векторы и матрицы — обычным начертанием в соответствии с ГОСТ 7.32
% (без полужирного, как принято в большинстве отечественных НИР)

% Окружение для алгоритмов
\newtheoremstyle{algodef}{8pt}{8pt}{\normalfont}{}{\bfseries}{.}{0.5em}{}
\theoremstyle{algodef}
\newtheorem{algorithm}{Алгоритм}[section]
\newtheorem{definition}{Определение}[section]
\newtheorem{statement}{Утверждение}[section]
